\ansclass{5}{Compute \& Services}
\ans{1} (a) \final{JWT or Redis}; today a fan is logged out when the load balancer picks another replica. (b) \final{Object store} (upload with a pre-signed URL); the poster exists on one replica only and vanishes with it. (c) \final{SQL row with a unique constraint} (or Redis \texttt{SET NX} with an expiry); two replicas can hold A-12 for two fans. (d) \final{Redis, or a short TTL}; replicas show different prices for a long time. (e) \final{A scheduler with a lock or a leader} (or a CronJob); it runs once per replica. (f) \final{Redis \texttt{INCR}} or the database; each replica counts only its share and the count resets on restart.
\careful{fixing it with sticky sessions: the state still dies with the replica, and the load becomes uneven.}
\ans{2} (a) $10 / 0.020 =$ \final{500 req/s}. (b) $4\,000/500 =$ \final{8 at 100\%} (zero headroom, latency explodes); $4\,000/(500 \times 0.65) = 12.3 \to$ \final{13}; with a spare \final{14}. (c) $500/20 =$ \final{25 replicas at most}; 100 replicas would open 2\,000 connections, more than PostgreSQL accepts, and the 20 ms of work per request stays: cost without gain. (d) $\lceil 6\,000/(500 \times 0.6) \rceil + 1 = 20 + 1 =$ \final{21}.
\careful{sizing at 100\% and forgetting the spare: the first replica that dies, or the first slow second, pushes the others past the knee.}
\ans{3} (a) $100 \times 0.8 =$ \final{80 workers}. (b) $40/0.8 =$ \final{50 jobs/s}; backlog grows $100 - 50 = 50$/s, so $50 \times 600 =$ \final{30\,000 jobs}; the last fan waits $30\,000/50 = 600$ s $=$ \final{10 min}. (c) Drains at $50 - 10 = 40$/s: $30\,000/40 =$ \final{750 s $\approx$ 12.5 min}. (d) 40: \final{no} (10 min). 60: capacity 75/s, backlog $25 \times 600 = 15\,000$, last wait $15\,000/75 = 200$ s $\approx$ \final{3.3 min: yes}, at three quarters of the cost of 80 workers. (e) Only 100 sends/s succeed; the rest collect \final{429s and retries}: the downstream limit, not the worker count, sets the real cap.
\careful{dividing the backlog by the arrival rate instead of the drain rate (or using the total arrivals as the backlog).}
\ans{4} (a) $2\,000 \times 0.2 =$ \final{400 environments}; $1\,500 \times 0.2 =$ \final{300}, i.e. up to 300 database connections: use a pooler (PgBouncer, RDS Proxy), a queue or reserved concurrency. (b) $0.20 + 0.5 \times 0.1 \times 0.0000166667 \times 10^6 \approx$ \final{US\$ 1.03 per million}. (c) $\lceil 6 \times 90/60 \rceil =$ \final{9}. (d) Metrics, the HPA loop, scheduling, image pull and warm-up add up to \final{1--4 minutes}, and a replica takes 90 s to be ready: the spike is over the target before help arrives. Pre-scale: $20\,000/(500 \times 0.6) = 66.7 \to$ \final{67 replicas at 09:45}, back down at 10:45 (and check the database can take them).
\careful{assuming elasticity is instant: reactive autoscaling follows a spike, it does not meet it.}
\ans{5} (a) \final{10 e-mails per fan}. Fixes: a Kubernetes \final{CronJob} with \texttt{concurrencyPolicy: Forbid}, a \final{distributed lock} (ShedLock) or leader election, plus an idempotent job (a ``reminder sent'' flag). (b) The lock alone does not: A believes it still holds it. The write carries the token and the storage \final{rejects 33 because it has seen 34} (fencing), or the job uses a conditional update. (c) Sweeper: simple, frees seats up to 30 s late; delayed message: on time, but one more component; \final{lazy expiry}: always correct, because a hold past \texttt{expires\_at} counts as free when read. With lazy expiry plus a sweeper for cleanup, a 5-minute scheduler outage \final{loses no seat}. (d) The opening time is data: \texttt{inventory} rejects holds before \texttt{onSaleAt} $=$ 10:00:00; the scheduler only prepares (warm-up at 09:45). The sale must not depend on a job firing on time.
\careful{trusting a lock for correctness. Locks expire; the data must defend itself.}
\ans{6} (a) \final{UNKNOWN}: the charge may have happened. Retry \final{with the same idempotency key} derived from the order (\texttt{ord\_9}), or query the provider; never a new key. (b) Refused: \final{retry}. Timeout after sending: \final{retry with the same key} (or query). 503: \final{retry} with backoff and jitter. 429: \final{wait for \texttt{Retry-After}}. 402 and 422: \final{do not retry}. (c) \final{Valid signature} (HMAC over id, timestamp and raw body), \final{timestamp within 300 s}, \final{event id not seen before} (dedupe), and a \final{monotonic transition} (a conditional \texttt{UPDATE \dots{} WHERE status IN (\dots)}: 0 rows means duplicate or stale); answer 2xx fast and process after. (d) \final{Reconciliation}: a sweeper every few minutes polls payments still UNKNOWN or PENDING after 10 min, and a daily job compares with the provider's settlement report. (e) Circuit breaker open, fail fast with a clear message, and \final{extend the holds by the outage}, so fans do not lose their seats.
\careful{treating a timeout as a failure and charging again with a new key.}
\ans{7} A good drawing: Client $\to$ API Gateway $\to$ \texttt{order} (App Server, SQL) \final{sync} $\to$ \texttt{inventory} hold \final{sync} $\to$ \texttt{payment} $\to$ Payment provider \final{sync with timeout and key}; the provider's webhook enters through a Serverless receiver; \texttt{OrderConfirmed} goes \final{async} through a Message Queue to Workers (PDF) and Notifications $\to$ Email; a Scheduler runs the hold sweeper and reconciliation. (a) \final{\texttt{inventory}}: its scaling profile (hot during on-sales) and strong consistency differ from the rest; extract it first, keep the others as modules. (b) $0.999^3 \approx$ \final{99.7\%}: every split on the synchronous path costs availability and latency, so split only where a force pays for it. (c) Serverless for \final{the webhook receiver and image processing} (spiky, short); \final{never the inventory} (hot path, cold starts, database connections).
\careful{splitting by technical layer or into entity services that share one database: a distributed monolith with all the costs and none of the gains.}
